\documentclass[11pt]{article}
\usepackage{cis3300-accessible}
\usepackage{amssymb}
% Keep every table in reading order, where it appears in the source.
\usepackage{float}
% Same as accessibletable in cis3300-accessible.sty, but pinned in place
% ([H]) so the grading tables stay in reading order instead of floating.
\renewenvironment{accessibletable}[3]{%
  \begin{table}[H]
  \centering
  \caption{#1}
  \label{#2}
  \begin{tabular}{#3}
}{%
  \end{tabular}
  \end{table}
}

% This document is mostly Java. Listings that are not Java (shell
% commands, CSV data, and expected output) override this locally
% with language={}.
\lstset{language=Java}

\AssignmentMeta{Task W7: Arcade High Score Tracker}{CIS 3300, Intermediate Object-Oriented Programming}

\begin{document}
\AssignmentTitleBlock
\AIDisclaimer

\section{Scenario}

You have been asked to help a small retro arcade turn a messy high score log
into a clean, ranked report. The arcade keeps every score from every machine
in one plain CSV file, in whatever order the scores happen to come in. Your
program reads that file, ranks every entry from the highest score to the
lowest, and writes a new CSV file with the finished report.

For this task, you will do all of the work from the Linux command line:
creating your project folders, writing your Java class in a plain text
editor, compiling it, and running it. You will not use Eclipse or any other
IDE for this task. The point is to see, by hand, everything an IDE normally
does for you behind the scenes.

This is an individual task. Complete it on your own, the same as any other
Task or Lab in this course.

\section{Learning Objectives}

By the end of this task you will have practiced:

\begin{itemize}
  \item laying out a Java project by hand from the command line, in the same
    \texttt{src}, \texttt{bin}, and \texttt{files} arrangement Eclipse has
    been building for you automatically all semester
  \item creating, editing, and saving a plain text file without an IDE
  \item compiling a Java source file with \texttt{javac} and running the
    compiled class with \texttt{java}, including the \texttt{-d} and
    \texttt{-cp} flags
  \item reading a CSV file line by line with \texttt{Scanner}, splitting each
    line into fields, and storing the results in an \texttt{ArrayList}
  \item sorting a list with a simple algorithm you already know, and writing
    formatted output to a new file with \texttt{PrintWriter}
  \item reading an existing program closely enough to make a small,
    deliberate change to it yourself
\end{itemize}

\section{About Your Screenshots}

Throughout this task, you will be asked to take a screenshot to show that
you completed a step. A few things apply to all of them.

\begin{itemize}
  \item Take one screenshot per marked step below. There are six in total,
    covering Sections 5 through 10.
  \item A screenshot does not need to show everything you did for that step.
    Get as much of the relevant terminal output on screen as you reasonably
    can.
  \item Each image must be an actual screenshot taken on your computer, not
    a photo taken with a phone or camera. If you do not already know how to
    take a screenshot on Xubuntu, look this up, for example a keyboard
    shortcut or a built in screenshot tool.
  \item \textbf{File naming.} Name each screenshot to match the section it
    documents: \texttt{taskw07\_5.png} for Section 5, \texttt{taskw07\_6.png}
    for Section 6, and so on through \texttt{taskw07\_10.png}.
\end{itemize}

\section{A Note on Editors}

Everything in this task happens in a plain text editor and the terminal,
not inside an IDE. Xubuntu already has \texttt{nano} and
\texttt{vim} installed, so both are ready to use right away.

\subsection{vim, the editor you are most likely to see me use}

Open a file with vim by typing \texttt{vim FILENAME}. Vim starts in
\textbf{Normal mode}, where the keys you press are commands rather than
text. Press \texttt{i} to switch to \textbf{Insert mode}, where vim behaves
like an ordinary editor and you can type normally. When you are done typing,
press \texttt{Esc} to return to Normal mode.

From Normal mode, a few commands that begin with a colon save the file and
quit.

\begin{lstlisting}[language={}]
:w    # save (write) the file
:q    # quit
:wq   # save and quit in one step
:q!   # quit without saving
\end{lstlisting}

\noindent\textbf{\color{coaNavy}If you ever feel stuck.} Press \texttt{Esc}
to make sure you are in Normal mode, then type \texttt{:q!} and press Enter.
That leaves vim without saving, no matter what state you were in. Almost
everyone needs this at least once, so it is worth remembering early.

\subsection{nano, a simpler fallback}

If you would rather start with something gentler than vim, \texttt{nano} is
a friendlier editor. Open a file with \texttt{nano FILENAME}, type normally,
then save with \texttt{Ctrl+O}, press Enter to confirm the file name, and
exit with \texttt{Ctrl+X}.

\subsection{GVim and NeoVim, two more options}

You are not required to install or use either of these for this task, but
you should know they exist.

\textbf{GVim} is a graphical version of vim. It opens in its own window with
menus and mouse support, but the underlying editor, and every command shown
above, is the same vim you already know. If you want to try it, install it
with:

\begin{lstlisting}[language={}]
sudo apt install vim-gtk3
\end{lstlisting}

\noindent then open a file with \texttt{gvim FILENAME}.

\textbf{NeoVim} is a modern, actively maintained reimplementation of vim. It
runs in the terminal the same way vim does, keeps the same modes and most of
the same commands, and adds a more current plugin system and configuration
language underneath. If you want to try it, install it with:

\begin{lstlisting}[language={}]
sudo apt install neovim
\end{lstlisting}

\noindent then open a file with \texttt{nvim FILENAME}.

\begin{accessibletable}{A quick comparison of the four editors mentioned in this task.}{tab:editors}{p{2.4cm}p{3.4cm}p{7.8cm}}
\toprule
\textbf{Editor} & \textbf{Interface} & \textbf{Notes} \\
\midrule
vim & Terminal only & Installed already. The editor you are most likely to see me use. \\
nano & Terminal only & Installed already. The simplest option to start with. \\
GVim & Graphical window & Same vim underneath, with menus and mouse support. Optional install. \\
NeoVim & Terminal, with optional graphical front ends & A modern, plugin friendly reimplementation of vim. Optional install. \\
\bottomrule
\end{accessibletable}

\noindent\textbf{\color{coaNavy}One more simple option.} Xubuntu also
includes \textbf{Mousepad}, a basic graphical text editor, if you would rather
work outside the terminal. Open it from the Applications menu, or type
\texttt{mousepad FILENAME} from inside your project folder. Mousepad uses its
own save dialog, so pay close attention to where each file is saved. If you
open it from the Applications menu, the dialog may default to your home
folder or Desktop. Make sure every file ends up in the correct place inside
your project folder, for example \texttt{files/scores.csv} or
\texttt{src/arcade/ArcadeScores.java}.

\subsection{These editors in the wider industry}

Vim and its descendants are not just classroom tools. Many professional
developers, particularly in backend, systems, and open source work, use vim
or NeoVim as their everyday editor, often alongside a terminal multiplexer
such as tmux. A few videos, found with web search, give a sense of how this
looks in practice.

\begin{itemize}
  \item \href{https://www.youtube.com/@Fireship}{Fireship (YouTube
    channel)}, a widely watched programming channel whose short "100 Seconds
    of Code" series includes quick explainers on both vim and NeoVim.
  \item \href{https://youtube.com/ThePrimeagen}{ThePrimeagen (YouTube
    channel)}, run by a former Netflix engineer who now codes and streams
    almost entirely in NeoVim, with videos that walk through a real,
    terminal based development setup.
  \item \href{https://www.youtube.com/watch?v=tNZnLkRBYA8}{Lex Fridman
    Podcast \#461: ThePrimeagen}, a long form conversation that includes a
    section on why he uses NeoVim and how a terminal based workflow fits
    into a professional engineering career.
\end{itemize}

\ResourcesDisclaimer

\section{Project Setup}

\noindent\textbf{\color{coaNavy}A note on your environment.} Your Xubuntu VM
already includes Firefox, so you can open this assignment, Moodle, and GitHub
entirely from inside the VM without switching back to your host machine. You
may complete this task on a different Linux machine instead of the VM if you
prefer, but you are responsible for installing and configuring everything
yourself (a JDK, git, and so on). If you need a refresher on using the VM,
see Lab W1.

Open a terminal (\textbf{Applications > Terminal Emulator}, or
\texttt{Ctrl+Alt+T}) and build the following folder structure by hand. This
mirrors the \texttt{src}, \texttt{bin}, and \texttt{files} layout Eclipse has
been building for you automatically all semester: \texttt{src} for your Java
source files, \texttt{bin} for the compiled class files, and \texttt{files}
for anything your program reads or writes.

\begin{enumerate}[itemsep=2pt]
  \item Create the project folders.
\begin{lstlisting}[language={}]
mkdir -p ~/cis3300/taskw07/src/arcade
mkdir -p ~/cis3300/taskw07/bin
mkdir -p ~/cis3300/taskw07/files
\end{lstlisting}
  \item Move into the project folder. You will run every command for the
    rest of this task from here.
\begin{lstlisting}[language={}]
cd ~/cis3300/taskw07
\end{lstlisting}
  \item Confirm the structure.
\begin{lstlisting}[language={}]
ls -R
\end{lstlisting}
\end{enumerate}

Take a screenshot of the \texttt{ls -R} output showing the \texttt{src},
\texttt{bin}, and \texttt{files} folders, and save it as
\texttt{taskw07\_5.png}.

\section{Creating the Data File}

The arcade's raw score log is a simple CSV file: a header row, then one row
per score, with the player's name, the game title, and the score as an
integer.

Open a new file at \texttt{files/scores.csv} in the editor of your choice
(vim, nano, GVim, NeoVim, or Mousepad) and type the following exactly as written,
including the header row.

\begin{lstlisting}[language={}]
name,game,score
Maria,Galaxy Invaders,15200
Deshawn,Retro Racer,9800
Priya,Galaxy Invaders,21300
Wyatt,Pixel Dungeon,7600
Alexis,Retro Racer,12500
Noah,Pixel Dungeon,15200
Grace,Galaxy Invaders,18400
Malik,Retro Racer,9800
Sofia,Pixel Dungeon,20100
Ethan,Galaxy Invaders,11000
\end{lstlisting}

It is fine to copy and paste this data instead of typing it, but check the
result carefully afterward. A paste from a PDF does not always come through
cleanly, and extra spaces can sneak in around the commas. There should be no
space before or after any comma in the file.

Save the file, then confirm its contents from the command line.

\begin{lstlisting}[language={}]
cat files/scores.csv
\end{lstlisting}

Take a screenshot of the \texttt{cat} output and save it as
\texttt{taskw07\_6.png}.

\section{Creating the Java Source File}

Below is the full program you will build. Read through it once before you
start typing, so you know roughly what each part is doing.

\texttt{ArcadeScores} does three things, in order: it reads every row of
\texttt{files/scores.csv} into a list, it sorts that list from the highest
score to the lowest using a simple selection sort, and it writes a new file,
\texttt{files/ranked\_scores.csv}, with a rank number added to every row.

\subsection{Reading the file}
The \texttt{readScores} method opens \texttt{files/scores.csv} with a
\texttt{Scanner}, reads and discards the header row, and then reads every
remaining line. Each line is split on its commas into a \texttt{String[]}
with three elements (name, game, score), and every one of those arrays is
added to an \texttt{ArrayList}.

\subsection{Sorting the list}
The \texttt{sortByScoreDescending} method is a selection sort. For each
position in the list, it scans the rest of the list for the highest
remaining score, then swaps that row into the current position. This is the
same idea as sorting a hand of cards by repeatedly picking the best card
left and moving it to the front.

\subsection{Writing the report}
The \texttt{writeRanking} method opens \texttt{files/ranked\_scores.csv}
with a \texttt{PrintWriter}, writes a header row, and then writes one line
per row in the now sorted list, numbering them starting at 1.

Create the file \texttt{src/arcade/ArcadeScores.java} in your editor and
type the following exactly as written.

\begin{lstlisting}
package arcade;

import java.io.File;
import java.io.FileNotFoundException;
import java.io.FileWriter;
import java.io.IOException;
import java.io.PrintWriter;
import java.util.ArrayList;
import java.util.List;
import java.util.Scanner;

/*
 * ArcadeScores reads a raw high score log from files/scores.csv and
 * writes a ranked report to files/ranked_scores.csv. Each input row
 * has three fields: a player name, a game title, and an integer
 * score. The output adds a rank column, with the highest score
 * ranked first.
 */
public class ArcadeScores {

    public static void main(String[] args) throws IOException {

        String inputPath = "files/scores.csv";
        String outputPath = "files/ranked_scores.csv";

        List<String[]> rows = readScores(inputPath);
        sortByScoreDescending(rows);
        writeRanking(rows, outputPath);

        System.out.println("Wrote " + rows.size() + " ranked scores to " + outputPath);
    }

    // Reads scores.csv, skips the header row, and returns each data
    // row as a String array: {name, game, score}.
    private static List<String[]> readScores(String path) throws FileNotFoundException {
        List<String[]> rows = new ArrayList<>();
        Scanner fileScanner = new Scanner(new File(path));

        if (fileScanner.hasNextLine()) {
            fileScanner.nextLine(); // skip the header row
        }

        while (fileScanner.hasNextLine()) {
            String line = fileScanner.nextLine();
            if (line.isBlank()) {
                continue;
            }
            rows.add(line.split(","));
        }

        fileScanner.close();
        return rows;
    }

    // Sorts rows by score, highest first, using a simple selection
    // sort. On each pass, it finds the highest remaining score and
    // swaps it into place.
    private static void sortByScoreDescending(List<String[]> rows) {
        for (int i = 0; i < rows.size() - 1; i++) {
            int maxIndex = i;
            for (int j = i + 1; j < rows.size(); j++) {
                int scoreJ = Integer.parseInt(rows.get(j)[2]);
                int scoreMax = Integer.parseInt(rows.get(maxIndex)[2]);
                if (scoreJ > scoreMax) {
                    maxIndex = j;
                }
            }
            String[] temp = rows.get(i);
            rows.set(i, rows.get(maxIndex));
            rows.set(maxIndex, temp);
        }
    }

    // Writes the ranked rows to a new CSV file, numbering them
    // starting at 1 in the order they are given.
    private static void writeRanking(List<String[]> rows, String path) throws IOException {
        PrintWriter writer = new PrintWriter(new FileWriter(path));
        writer.println("rank,name,game,score");

        int rank = 1;
        for (String[] row : rows) {
            String name = row[0];
            String game = row[1];
            String score = row[2];
            writer.println(rank + "," + name + "," + game + "," + score);
            rank++;
        }

        writer.close();
    }
}
\end{lstlisting}

It is fine to copy and paste this code instead of typing it, but you will
likely need to fix the line spacing and tabs afterward, since a paste from a
PDF does not preserve indentation cleanly.
\textbf{\color{coaNavy}Formatting matters.} Code that is not formatted nicely
will not receive credit.

Save the file when you are done typing.

Take a screenshot of your editor with \texttt{ArcadeScores.java} open,
showing at least part of the code on screen, and save it as
\texttt{taskw07\_7.png}.

\section{Compiling and Running}

The \texttt{javac} command compiles your source file into bytecode, and the
\texttt{java} command runs it. Because your class lives in the
\texttt{arcade} package, both commands need a little more than the bare file
name.

\begin{enumerate}[itemsep=2pt]
  \item Compile the source file into the \texttt{bin} folder.
\begin{lstlisting}[language={}]
javac -d bin src/arcade/ArcadeScores.java
\end{lstlisting}
    The \texttt{-d bin} flag tells \texttt{javac} where to put the compiled
    \texttt{.class} file. Because the class declares
    \texttt{package arcade;}, \texttt{javac} creates an \texttt{arcade}
    folder inside \texttt{bin} for you.
  \item Confirm the class file was created.
\begin{lstlisting}[language={}]
ls bin/arcade
\end{lstlisting}
    You should see \texttt{ArcadeScores.class}. This is the compiled
    bytecode; it is what \texttt{java} actually runs, not the \texttt{.java}
    source file.
  \item Run the compiled class.
\begin{lstlisting}[language={}]
java -cp bin arcade.ArcadeScores
\end{lstlisting}
    The \texttt{-cp bin} flag tells \texttt{java} where to look for compiled
    classes, and \texttt{arcade.ArcadeScores} is the fully qualified class
    name (package name, a dot, then the class name), not a file name.
\end{enumerate}

If everything worked, you will see a line similar to this.

\begin{lstlisting}[language={}]
Wrote 10 ranked scores to files/ranked_scores.csv
\end{lstlisting}

Take one screenshot showing the compile command, the \texttt{ls bin/arcade}
output, the run command, and its printed output together in one screen, and
save it as \texttt{taskw07\_8.png}.

\section{Checking the Output}

Look at the file your program wrote.

\begin{lstlisting}[language={}]
cat files/ranked_scores.csv
\end{lstlisting}

If your program worked correctly and you typed the data file exactly as
given, your output should match the following exactly.

\begin{lstlisting}[language={}]
rank,name,game,score
1,Priya,Galaxy Invaders,21300
2,Sofia,Pixel Dungeon,20100
3,Grace,Galaxy Invaders,18400
4,Noah,Pixel Dungeon,15200
5,Maria,Galaxy Invaders,15200
6,Alexis,Retro Racer,12500
7,Ethan,Galaxy Invaders,11000
8,Malik,Retro Racer,9800
9,Deshawn,Retro Racer,9800
10,Wyatt,Pixel Dungeon,7600
\end{lstlisting}

If your file does not match, do not move on yet. Compare your
\texttt{scores.csv} against Section 6 character for character, and compare
your \texttt{ArcadeScores.java} against Section 7 the same way. A single
missing comma or a mistyped variable name is the most common cause.

Take a screenshot of your own \texttt{cat} output and save it as
\texttt{taskw07\_9.png}.

\section{The Challenge: Add a Medal Column}

The arcade owner has one more request. The ranked report should show a
medal for the top three scores: \texttt{Gold} for rank 1, \texttt{Silver}
for rank 2, and \texttt{Bronze} for rank 3. Every other row should show a
dash, \texttt{-}, in that column instead.

The finished output should look like this, with your own column added.

\begin{lstlisting}[language={}]
rank,name,game,score,medal
1,Priya,Galaxy Invaders,21300,Gold
2,Sofia,Pixel Dungeon,20100,Silver
3,Grace,Galaxy Invaders,18400,Bronze
4,Noah,Pixel Dungeon,15200,-
\end{lstlisting}

You are not given the code for this part.

\begin{itemize}
  \item Only the part of the program that writes the output needs to
    change. You do not need to touch how the file is read or how it is
    sorted.
  \item The header row needs a new column name.
  \item Inside the loop that writes each row, you already have the current
    \texttt{rank} as a number. Decide what that row's medal should be based
    on that number, and add it to the line you write.
\end{itemize}

Edit \texttt{src/arcade/ArcadeScores.java} in your editor, save it, then
recompile and rerun it the same way you did in Section 8.

\begin{lstlisting}[language={}]
javac -d bin src/arcade/ArcadeScores.java
java -cp bin arcade.ArcadeScores
\end{lstlisting}

Check the result.

\begin{lstlisting}[language={}]
cat files/ranked_scores.csv
\end{lstlisting}

Take one screenshot showing your recompile, rerun, and the final
\texttt{cat} output together in one screen, and save it as
\texttt{taskw07\_10.png}.

\section{Grading}

This task is worth 29 points. Each item below is either met or not met; there is no partial credit within an item. Each table lists one category of items with its point value.

\medskip
\noindent\textbf{\color{coaNavy}Formatting.} Severely poor formatting may result in a zero for the entire task.

\begin{accessibletable}{Grading checklist, Submission and Setup (1 point each, 4 points).}{tab:g-setup}{c L{11.6cm} c}
\toprule
\textbf{Met} & \textbf{Item} & \textbf{Points} \\
\midrule
$\square$ & \texttt{taskw07} folder is present in the \texttt{cis3300} repository, following the structure shown in Section 5 & 1 \\
$\square$ & \texttt{ArcadeScores.java} compiles without errors from the command line & 1 \\
$\square$ & \texttt{ArcadeScores.java} runs without errors from the command line & 1 \\
$\square$ & \texttt{bin} folder is not included in the submitted repository & 1 \\
\bottomrule
\end{accessibletable}

\begin{accessibletable}{Grading checklist, Screenshots (1 point each, 6 points).}{tab:g-shots}{c L{11.6cm} c}
\toprule
\textbf{Met} & \textbf{Item} & \textbf{Points} \\
\midrule
$\square$ & \texttt{taskw07\_5.png} is present, showing the \texttt{ls -R} output from Section 5 & 1 \\
$\square$ & \texttt{taskw07\_6.png} is present, showing the \texttt{cat} output from Section 6 & 1 \\
$\square$ & \texttt{taskw07\_7.png} is present, showing the editor open with \texttt{ArcadeScores.java} from Section 7 & 1 \\
$\square$ & \texttt{taskw07\_8.png} is present, showing the compile and run output from Section 8 & 1 \\
$\square$ & \texttt{taskw07\_9.png} is present, showing the \texttt{cat} output from Section 9 & 1 \\
$\square$ & \texttt{taskw07\_10.png} is present, showing the recompile, rerun, and final \texttt{cat} output from Section 10 & 1 \\
\bottomrule
\end{accessibletable}

\begin{accessibletable}{Grading checklist, Base Program Correctness (2 points each, 10 points).}{tab:g-base}{c L{11.6cm} c}
\toprule
\textbf{Met} & \textbf{Item} & \textbf{Points} \\
\midrule
$\square$ & \texttt{files/scores.csv} contains the data exactly as given, unmodified & 2 \\
$\square$ & \texttt{readScores} correctly reads all ten rows, skipping the header row & 2 \\
$\square$ & \texttt{sortByScoreDescending} correctly sorts the list from highest score to lowest & 2 \\
$\square$ & \texttt{writeRanking} writes \texttt{files/ranked\_scores.csv} with a header row and rank numbers starting at 1 & 2 \\
$\square$ & Output in \texttt{files/ranked\_scores.csv} matched the Expected Output in Section 9 before the challenge change & 2 \\
\bottomrule
\end{accessibletable}

\begin{accessibletable}{Grading checklist, The Challenge: Medal Column (2 points each, 4 points).}{tab:g-challenge}{c L{11.6cm} c}
\toprule
\textbf{Met} & \textbf{Item} & \textbf{Points} \\
\midrule
$\square$ & Header row includes a medal column, and the reading and sorting logic is unchanged from the base program & 2 \\
$\square$ & Every row shows the correct medal (\texttt{Gold}, \texttt{Silver}, \texttt{Bronze}, or a dash) for its rank & 2 \\
\bottomrule
\end{accessibletable}

\begin{accessibletable}{Grading checklist, Code Quality (1 point each, 5 points).}{tab:g-quality}{c L{11.6cm} c}
\toprule
\textbf{Met} & \textbf{Item} & \textbf{Points} \\
\midrule
$\square$ & Indentation is consistent throughout \texttt{ArcadeScores.java} & 1 \\
$\square$ & Code is well formatted and follows standard Java coding conventions & 1 \\
$\square$ & The header comment block is still present at the top of the file & 1 \\
$\square$ & No leftover debug output or commented-out code beyond what was provided & 1 \\
$\square$ & No unrequested extra code or output added beyond what the assignment asked for & 1 \\
\bottomrule
\end{accessibletable}

\section{Submission}

Remember that your Xubuntu VM includes Firefox, so you can submit your work
directly from the VM without moving files to your host machine first.

\begin{enumerate}[itemsep=2pt]
  \item Inside your \texttt{cis3300} repository, you should have a
    \texttt{taskw07} folder containing your \texttt{src} and \texttt{files}
    folders, and your six screenshots. Do not include the \texttt{bin}
    folder; compiled class files do not belong in your repository.
  \item Add the \texttt{taskw07} folder to your repository, using either
    option:
    \begin{itemize}
      \item commit and push from the terminal with git, or
      \item upload the folder through the GitHub website in Firefox (leave
        out the \texttt{bin} folder).
    \end{itemize}
  \item Return to Task W7 in Moodle and type \texttt{done} in the text
    submission box. Do not upload images to the submission box.
\end{enumerate}

\end{document}
